\documentclass[twocolumn,10pt,journal,compsoc]{IEEEtran}

\usepackage{cite}
\usepackage{amsmath,amssymb,amsfonts}
\usepackage{graphicx}
\usepackage{textcomp}
\usepackage{xcolor}
\usepackage{hyperref}

\begin{document}

\title{Physics-Informed DeepONets with Residual-Based Adaptive Refinement for 2D Solid Mechanics}

\author{Amith Gowda \\
\textit{Department of Scientific Computing}
}

\maketitle

\begin{abstract}
This report investigates the application of Physics-Informed Deep Operator Networks (PI-DeepONets) combined with Residual-based Adaptive Refinement (RAR) to solve 2D linear elasticity problems. We focus on the classic Kirsch problem—a flat plate containing a circular hole under uniaxial tension. Through four experimental arms (uniform baseline, collocation-only RAR, load-only RAR, and combined RAR), we evaluate the efficiency and accuracy of adaptive training strategies in operator learning for solid mechanics.
\end{abstract}

\begin{IEEEkeywords}
Physics-Informed Neural Networks, DeepONet, Operator Learning, Residual-based Adaptive Refinement, Solid Mechanics, Kirsch Problem.
\end{IEEEkeywords}

\section{Introduction}
Solving partial differential equations (PDEs) representing physical systems is core to engineering. Traditional methods like the Finite Element Method (FEM) are computationally expensive for parametric studies. Operator learning frameworks, such as DeepONets \cite{lu2021learning}, learn the mapping between function spaces (e.g., boundary conditions to displacement fields). By adding physics constraints, Physics-Informed DeepONets \cite{wang2021learning} train without labeled data.

In this work, we implement Residual-based Adaptive Refinement (RAR) to actively add collocation points in areas showing high PDE residuals and update loading function training sets.

\section{Methodology}
\subsection{PI-DeepONet Architecture}
A DeepONet consists of a branch network that processes the input function $T(y)$ evaluated at sensor locations, and a trunk network that processes coordinates $(x,y)$. The outputs are combined via dot products to yield displacements $u(x,y)$ and $v(x,y)$.

\subsection{2D Linear Elasticity & Navier-Cauchy Equations}
Under plane stress conditions, the equilibrium equations are:
\begin{equation}
\frac{\partial \sigma_{xx}}{\partial x} + \frac{\partial \sigma_{xy}}{\partial y} = 0
\end{equation}
\begin{equation}
\frac{\partial \sigma_{xy}}{\partial x} + \frac{\partial \sigma_{yy}}{\partial y} = 0
\end{equation}
where stresses are related to strains through Hooke's Law.

\subsection{Residual-Based Adaptive Refinement (RAR)}
Instead of uniform sampling, RAR evaluates the PDE residuals on a large candidate pool of coordinate locations or loading functions. The samples yielding the highest residuals are iteratively added to the training batch during optimization.

\section{Experimental Setup}
We study four configurations:
\begin{itemize}
    \item \textbf{Baseline:} Uniform sampling of coordinates and loading functions.
    \item \textbf{Collocation RAR:} Adaptive addition of spatial coordinates.
    \item \textbf{Load RAR:} Adaptive addition of boundary load functions.
    \item \textbf{Combined RAR:} Combined adaptive refinement of both spatial coordinates and load functions.
\end{itemize}

\section{Results \& Discussion}
(Detailed results comparing L2 errors and Stress Concentration Factors (SCF) go here.)

\section{Conclusion}
We demonstrated that combining collocation-based and load-based RAR significantly improves the convergence and boundary prediction accuracy of PI-DeepONet models for solid mechanics problems.

\bibliographystyle{IEEEtran}
\bibliography{references}

\end{document}
